\chapter{有监督微调}
\label{chap:nlp-supervised-finetuning}

用户把活动通知交给模型，要求“用一句话说明外部人员怎样参加”。一个能续写通知的预训练模型，可能继续补写主办单位，也可能复述整篇材料。即使它知道日期和报名方式，也还需要学会把用户要求当作当前目标，把通知当作回答依据，并按指定方式结束回答。第\ref{chap:nlp-large-language-models}章解释了模型怎样从语料中学习条件概率；本章研究怎样利用示范，把已有模型调整为能够完成交互要求的模型。

这里的\term{有监督微调}（supervised fine-tuning，SFT），是从已有参数出发，利用输入及期望回答，或者教师提供的概率分布，继续调整模型。它不必改变模型结构，也不意味着所有参数都要更新。为了判断一次微调究竟教会了什么，我们需要连起来考察四个环节：示范是否覆盖目标要求，内容怎样变成监督位置和损失，哪些参数按照什么强度更新，以及更新后能否完成未见过的要求。这四个环节相互制约；精选答案如果配错了历史，或者更新过强，仍可能产生不合用的模型。

\section{示范材料准备}
\label{sec:nlp-sft-demonstrations}

训练材料需要回答的是“在什么信息条件下，什么回答值得模仿”。这一问题先于词元计分和优化器选择。把所有优质文章放在一起，并不能自动得到优质交互示范，因为文章未必说明模型面对的要求，也未必保留作答时可见的依据。

\subsection{交互要求与示范覆盖}
\label{sec:nlp-sft-coverage}

仍以活动通知为例，同一份材料可以对应三种示范。用户问“外部人员怎么参加”时，回答需要同时包含报名条件和截止时间；用户问“地点从哪里改到哪里”，而材料只给出新地点时，回答应说明已知地点，避免编造原地点；用户要求查询不在材料中的个人报名信息时，回答需要说明信息不足，并提出合适的后续询问。这些回答分别示范正常作答、信息不足时的澄清和行为边界。它们对应不同的信息条件，不能只用一种笼统的拒绝模板替代。

从这些差别出发，可以把示范范围写成可检查的要求：输入可能包含哪些语言和材料，用户可能怎样表达目标，哪些格式或长度必须遵守，多轮历史中哪些限制需要延续，以及哪些请求应澄清或拒绝。覆盖清单的单位应当是会改变响应方式的条件。例如“询问日期”和“查询日期”只是近义表达，而“只依据通知回答”和“允许结合外部常识解释”改变了允许使用的信息，两者都值得保留。

多任务指令微调把不同任务改写成自然语言要求，让模型在共同的输入输出形式下学习响应。FLAN系列的规模化研究考察了任务种类、模型规模和推理示范等因素，提供了指令多样性能够帮助未见任务泛化的经验证据\scite{nlp-llm-flan}。这类结果支持扩大有意义的条件覆盖，不能推出“任务名称越多，任何模型都越好”。若新增材料只换措辞而不改变条件，或者答案彼此矛盾，数量增长可能没有带来所需信息。

\subsection{人工示范与合成扩展}
\label{sec:nlp-sft-acquisition}

确定需要覆盖的条件之后，才有依据选择示范来源。人工撰写可以直接控制答案依据、表达方式和例外情形，但需要相应的专业知识与复核时间。已有问答、标注任务和人工修订记录也能转为示范，前提是保留原始输入，确认答案没有使用模型在作答时看不到的材料。否则，训练会要求模型从缺失信息中猜出答案。

教师模型能够扩展种子示范。Self-Instruct从少量人工任务出发，生成新指令、输入和输出，再过滤生成材料\scite{wang2023selfinstruct}。这一循环增加了候选材料的范围，却没有使教师输出自动成为事实依据。例如，教师把报名截止时间从通知中的“16日18时”改为“活动开始前”，语句仍然自然，但已经改变了要求。合成扩展后仍须回到来源核验答案，并检查新指令是否真正增加覆盖。

对不合适的回答，还可以先提出批评，再生成修订答案。Constitutional AI的监督阶段利用明确原则引导模型批评并修改原回答，随后把修订结果用于监督学习\scite{bai2022constitutional}。在本章中，学习对象是修订后的示范；若原则用于比较两个回答并形成偏好标签，则属于第\ref{chap:nlp-preference-alignment}章的反馈准备。两种做法可以使用相同原则，但对应的监督对象不同。

教师还会把自己的偏差复制到许多样本中。若同一教师总用冗长开场，学生可能把开场长度误当作答案质量；若教师在某一领域知识不足，批量生成会把局部错误扩成大量一致答案。因此，合成材料应保留教师版本、生成条件和来源标记，并用人工样本或可独立核验的依据检查内容。混合多个教师能够增加差异，但不能代替这种核验。

\subsection{筛选、去重与集合划分}
\label{sec:nlp-sft-filtering}

一个候选示范能否保留，需要同时检查输入和回答。输入必须说明目标与必要条件；回答必须遵守要求、使用允许的信息，并且在语义、格式和行为边界上合适。遇到同一输入对应冲突答案时，应先查明是答案错误，还是输入遗漏了条件。例如，两个示范分别要求外部人员报名和直接到场，若差别来自人员身份，就应把身份写入输入，而不是让模型学习一个无条件的多数答案。

重复样本会改变实际训练权重。假设一份含有错误截止时间的答案被复制了20次，另一份正确答案只出现1次，按样本平均的训练目标会更重地计入错误版本。近重复也有类似影响，因此，去重既要看字符串，也要看指令、依据和答案是否表达同一训练事实。保留重复有时是有意加权，但应当作为取样决策记录下来。

LIMA用少量精心选择的示范微调已有的大模型，显示在其模型和检查设置下，精选数据也能形成多种交互行为\scite{nlp-llm-lima}。这一观察不能说明所有新能力都能由少量样本获得，也不能说明预训练模型缺失的领域知识会自动出现。它提醒我们把数据质量、已有能力和覆盖范围放在一起判断，而不是只比较示范条数。

训练、开发和最终检查材料应按可能泄漏的来源划分。同一通知的近义问法如果分散在三个集合中，检查结果主要反映对已见内容的重述；要检验新材料上的泛化，应把整份通知及其派生样本放在同一侧。这样得到的示范集合不仅包含输入和答案，还保留来源、条件、质量判断及划分依据，下一节再把它转换成模型能够计算的监督。

\section{监督样本与损失构造}
\label{sec:nlp-sft-objectives}

示范确定了应当模仿的回答，还没有确定模型看到怎样的序列，以及哪些位置产生梯度。把同一份对话按不同角色、历史和损失掩码组织，可以形成不同的学习问题。本节明确这些选择，使数据中的交互含义与计算目标一致。

\subsection{对话角色与可见历史}
\label{sec:nlp-sft-context}

对于一轮交互，记当前可见的要求、材料和历史为$x$，期望回答为$y=(y_1,\ldots,y_m)$。输入序列需要用模型约定的角色和分隔标记表示不同来源，例如全局要求、用户问题和模型回答。这些标记本身也会被编码为词元。它们能够帮助模型学习“谁说了什么”，却不能凭标记形式保证模型始终正确处理相互冲突的要求。

材料中的文字与当前有效要求需要在样本中区分。假设通知引用了一句“请忽略报名截止时间”，而用户要求忠实提取通知中的报名安排；这句引文只是被分析的内容。若数据把引文也当成新的交互命令，学生便可能学会服从材料内部的指示。需要通过清楚的材料边界和相应示范，教会模型在当前任务中怎样使用这些文字。后续仍须用未见的材料检查这种行为。

多轮样本还要确定历史可见范围。训练当前回答时，可以保留之前的用户与模型消息，令模型学习延续已有约束；如果只保留最后一句“那外部人员呢”，输入就丢失了所指活动。序列超过长度预算时，截断应保留影响答案的条件，或者丢弃无法恢复条件的样本。任意截断后仍沿用完整历史中的答案，会产生输入与监督不相容的问题。

\subsection{回答位置与示范损失}
\label{sec:nlp-sft-masked-loss}

输入输出组织好以后，可以利用真实回答的前缀逐词计算概率。例如预测“16日18时前报名”中的第二个词元时，模型看到的是示范中的第一个词元，而不是自己抽样得到的词元。这种以正确前缀提供条件的计算称为\term{教师强制}（teacher forcing）。它使整段示范能够提供监督，但训练所见前缀与模型自行生成的前缀仍可能不同。

为了只在需要学习的回答位置计分，令第$i$个编码序列为
\[
 z_i=(z_{i,1},\ldots,z_{i,T_i}),
\]
并令$a_{i,t}\in\{0,1\}$表示第$t$个位置是否接受监督。样本权重$w_i\geq0$表示这份示范在目标中的相对贡献。一个按监督词元平均的损失为
\begin{equation}
\label{eq:nlp-sft-masked-loss}
\begin{aligned}
 L_{\mathrm{SFT}}(\bm{\theta})
 &=-\frac{1}{Z}\sum_{i=1}^{n}w_i
   \sum_{t=1}^{T_i}a_{i,t}
   \log p_{\bm{\theta}}(z_{i,t}\mid z_{i,<t}),\\
 Z&=\sum_{i=1}^{n}w_i\sum_{t=1}^{T_i}a_{i,t}>0.
\end{aligned}
\end{equation}
这个归一化使损失反映每个被计分词元的平均负对数概率。若先对每份回答内部平均，再对样本平均，长回答与短回答的相对权重就会改变；两种写法应依据训练目的选择。

全序列计分令指令、材料和回答都参与预测；回答位置计分只令模型学习给定输入后的响应；多轮选择计分还可以只监督某些质量可靠的模型轮次。以“用户：日期？模型：18日”为例，回答位置掩码在“18日”及需要学习的结束标记处为1，在用户问题处为0。掩码为0只意味着该位置没有直接损失，不意味着其内容不参与前向计算，也不意味着前面位置的表示不会通过回答损失获得梯度。

式\eqref{eq:nlp-sft-masked-loss}在每个被计分位置上，都鼓励提高给定前缀下示范词元的概率；共享参数使不同位置的要求需要共同折中。它没有直接声明其他答案都错误。例如“18日举行”和“活动在18日举行”可能同样合适，单一硬答案却只沿其中一条词元路径提供监督。如果希望保留多个可接受延续之间的概率关系，就需要更丰富的教师信号。

\subsection{教师分布与蒸馏损失}
\label{sec:nlp-sft-soft-distillation}

教师可以在每个前缀$h$上给出整个词表的概率$q(v\mid h)$，学生给出$p_{\bm{\theta}}(v\mid h)$。与只记录一个目标词元相比，这种软监督保留了教师对不同延续的相对判断。它要求教师与学生使用可比较的预测单位；词表或分词方式不同，不能直接把同一坐标的概率相减，必须先建立单位之间的对应关系，或者改用生成序列监督。

为了使学生覆盖教师认为可能的词元，可以最小化前向KL散度；若希望学生较少向教师低概率词元分配质量，可以考察反向KL散度。对于给定前缀，两者分别为
\begin{equation}
\label{eq:nlp-sft-kl-directions}
\begin{aligned}
 D_{\mathrm{KL}}(q\Vert p_{\bm{\theta}})
 &=\sum_v q(v\mid h)\log\frac{q(v\mid h)}{p_{\bm{\theta}}(v\mid h)},\\
 D_{\mathrm{KL}}(p_{\bm{\theta}}\Vert q)
 &=\sum_v p_{\bm{\theta}}(v\mid h)\log\frac{p_{\bm{\theta}}(v\mid h)}{q(v\mid h)}.
\end{aligned}
\end{equation}
求和覆盖共同词表，出现正分子和零分母时，相应散度为无穷。教师固定时，第一式中$\sum_v q\log q$与学生参数无关，因此最小化它等价于最小化软标签交叉熵。第二式还包含学生自身的熵项，两种更新一般不同。

例如教师在两个可接受词元上给出$(0.6,0.4)$。只保留第一个词元的硬答案相当于使用$(1,0)$监督，已经删除了第二个词元也可接受的信息。在学生容量受限、无法同时拟合所有上下文时，不同散度还会影响它怎样分配这种误差。GKD把散度选择与训练前缀来源放在同一框架中考察，后者在第\ref{sec:nlp-sft-on-policy-distillation}节落实\scite{nlp-llm-gkd}。

温度也会改变监督分布。若教师词元分数为$b_v$，使用温度$\tau>0$后，概率正比于$\exp(b_v/\tau)$；较高温度使分布较平缓，但不能据此认定教师更不确定的答案更可信。作为进阶方向，熵感知在线蒸馏根据教师分布的不确定程度组合散度，讨论如何减少单纯反向KL对多种合理延续的压缩\scite{nlp-llm-entropy-opd}。这类选择仍需结合任务、学生容量和生成质量检查，不能由某一散度的名称决定。

\section{参数适配与微调更新}
\label{sec:nlp-sft-updates}

前一节已经给出损失。接下来要明确损失能够改变哪些变量，以及训练样本怎样进入更新。参数更新范围、数据配比和学习强度都会限制模型最终能够接近哪一种行为，不能把它们当作与监督内容无关的实现细节。

\subsection{模型起点与更新范围}
\label{sec:nlp-sft-parameter-scope}

设预训练或继续预训练得到的参数为$\bm{\theta}_0$。全参数微调以它为起点，允许损失更新模型中指定的全部权重；部分冻结只更新选中的层或参数；参数适配则保留原权重，并引入较少的新变量。无论采用哪一种方式，都需要确认输入分词、位置处理和对话标记与已有模型兼容。新增词元还需要相应的表示和输出参数，不能只在文本模板中写出一个模型从未见过的特殊符号。

把可训练变量记为$\bm{\phi}$，将实际模型参数写成$\bm{\theta}(\bm{\phi};\bm{\theta}_0)$，便得到统一的学习问题：在选定更新空间内最小化$L_{\mathrm{SFT}}(\bm{\theta}(\bm{\phi};\bm{\theta}_0))$。全参数更新令$\bm{\phi}$包含全部权重；冻结与适配改变的是这组变量及其映射。一次梯度更新可写为$\bm{\phi}^{(k+1)}=\bm{\phi}^{(k)}-\eta_k\nabla_{\bm{\phi}}L_{\mathrm{SFT}}$，其中$\eta_k>0$为步长，实际梯度通常由本轮小批次估计。梯度可以穿过冻结的计算到达上游可训练变量，冻结只表示相应权重本身不被优化器更新。

相较于同一基座的全参数更新，冻结参数或限制增量会缩小允许的权重变化范围，但这不直接等同于更好的泛化。若新行为只需要较小调整，受限空间可能已经足够；若模型缺少必要表示，过强限制也可能妨碍学习。因此，更新范围应当与目标差距和保持要求一起检查，下一小节给出可计算的受限形式。

\subsection{受限参数空间中的适配}
\label{sec:nlp-sft-low-rank}

考虑一个线性变换$\bm{W}_0\bm{h}$，其中$\bm{W}_0\in\mathbb{R}^{d_{\mathrm{out}}\times d_{\mathrm{in}}}$。若只需要学习权重变化，可以冻结$\bm{W}_0$，用两个较窄的矩阵构造增量：
\begin{equation}
\label{eq:nlp-sft-lora}
\begin{aligned}
 \bm{W}&=\bm{W}_0+s\bm{B}\bm{A},\\
 \bm{A}&\in\mathbb{R}^{r\times d_{\mathrm{in}}},\qquad
 \bm{B}\in\mathbb{R}^{d_{\mathrm{out}}\times r}.
\end{aligned}
\end{equation}
这里$r$是所选秩上界，$s$是缩放系数，常按$\alpha/r$设置，其中$\alpha>0$是指定的缩放超参数。LoRA由损失更新$\bm{A}$、$\bm{B}$，使增量秩至多为$r$\scite{ket-hu2021lora}。例如一个$4096\times4096$矩阵有约1678万个元素，取$r=8$时，两组适配参数共65536个；减少的是该位置的可训练变量，模型的原权重仍然存在。

初始化还决定适配从哪个函数出发。常用做法给$\bm A$非零随机初值、令$\bm B=0$，于是初始增量为零，模型仍从$\bm W_0$开始。若将两者都初始化为零，只有依赖$\bm B\bm A$的数据拟合损失时，两组梯度也同时为零：令$\bm G=\nabla_{\bm W}L$，有$\nabla_{\bm B}L=s\bm G\bm A^{\mathsf T}$、$\nabla_{\bm A}L=s\bm B^{\mathsf T}\bm G$，这样的初始化无法启动适配。在$\bm A$非零、$\bm B=0$时，$\bm B$可以先获得非零梯度；$\bm B$发生变化后，$\bm A$也能随之更新。

低秩约束也有表达代价。若所需增量是二维单位矩阵，它的秩为2，单个秩1增量就不能精确表示。把适配器放入多个层后，网络整体变化可以更复杂，但不能据此把每个局部低秩约束说成没有限制。适配位置、秩和更新强度都应通过目标任务与保持任务共同选择。

QLoRA进一步把冻结的基座权重量化保存，在计算时使用其重建值，并把梯度传给低秩适配器\scite{nlp-llm-qlora}。若量化和重建得到$\widetilde{\bm{W}}_0$，实际使用的是$\widetilde{\bm{W}}_0+s\bm{B}\bm{A}$。适配器学习可以缓解某些误差，但并不保证恢复原始基座，也不保证与全精度LoRA相同；这里被量化冻结的权重与被训练的适配参数必须分开理解。

另一种限制方式是把权重的幅度和方向分开。以按列归一化为例，DoRA可写成
\begin{equation}
\label{eq:nlp-sft-dora}
 \bm{W}_{:,j}=m_j\,
 \frac{\bm{V}_{:,j}}{\lVert\bm{V}_{:,j}\rVert_2},
 \qquad \bm{V}=\bm{W}_0+s\bm{B}\bm{A}.
\end{equation}
$m_j$为第$j$列的可学习幅度，分母要求该列非零；初始幅度取原权重列范数，并令$\bm B\bm A=0$，才会在这些非零列上恢复原权重。这样既能沿低秩增量改变方向，又能直接调整幅度\scite{nlp-llm-dora}。公式采用的列约定随权重存储方向而变，核心是显式分开幅度与归一化方向。它改变了可训练参数形式，收益仍取决于具体任务与设置。

\subsection{数据混合与更新强度}
\label{sec:nlp-sft-mixture-strength}

假设示范来自若干集合$\mathcal{D}_k$，分别覆盖信息提取、多轮澄清、结构化输出和行为边界。按概率$\rho_k$选择来源、再在来源内取样，等价于在相应取样条件下优化$\sum_k\rho_k\mathbb{E}_{\mathcal{D}_k}[\ell]$。若每份样本还使用额外权重，权重会与取样概率共同决定有效目标。不能一面提高某类的取样率，一面把观察到的行为变化全部归因于优化器。

更新强度由学习率、更新次数、样本重复和可训练空间共同决定。微调较久可能继续降低示范损失，却也可能使模型过度模仿特定表述，或者改变原有知识的使用方式。一个合适的检查过程是在若干检查点上，同时记录未见指令的完成情况和需要保持的能力，再根据预先确定的要求选择参数，而不是只选择训练损失最小的时刻。

TÜLU 3把数据整理、SFT、偏好优化和可验证奖励训练组成后训练流程，能够用来观察各环节的接口\scite{nlp-llm-tulu3}。本节只回用其中SFT的材料混合与学习安排。报告中后续偏好或奖励训练带来的变化，不应记作示范损失本身的效果；若要判断某种混合策略的贡献，需要在相同初始模型和其他条件下比较相应阶段。

\subsection{学生生成与在线蒸馏}
\label{sec:nlp-sft-on-policy-distillation}

固定示范前缀训练后，学生在实际生成时可能先走到一个示范没有出现的前缀。比如目标回答是“16日18时前报名”，学生却先生成“17日”，后续概率便是在错误日期条件下计算。继续只在正确示范前缀上监督，不能直接告诉学生在自己常到达的状态中应怎样分配下一步概率。

一个做法是让当前学生先生成回答，再在这些前缀上查询教师分布。令$\bm{\theta}^{(k)}$表示本轮采样用的学生参数，$d_{\mathrm{fixed}}$表示固定材料中的前缀分布，$d_{\bm{\theta}^{(k)}}$表示学生生成得到的前缀分布。将两种来源混合，可写成
\begin{equation}
\label{eq:nlp-sft-gkd}
\begin{aligned}
 \ell_D(h;\bm{\theta})
 &=D(q(\cdot\mid h),p_{\bm{\theta}}(\cdot\mid h)),\\
 L^{(k)}(\bm{\theta})
 &=(1-\lambda)\mathbb{E}_{h\sim d_{\mathrm{fixed}}}\ell_D(h;\bm{\theta})\\
 &\quad+\lambda\mathbb{E}_{h\sim d_{\bm{\theta}^{(k)}}}\ell_D(h;\bm{\theta}).
\end{aligned}
\end{equation}
其中$0\leq\lambda\leq1$控制学生前缀的比例，$D$是已选定的分布差异。更新时把本轮采到的前缀视为固定数据，梯度作用于分布匹配损失，不穿过离散采样过程；更新后可以重新生成下一轮前缀。这正是GKD所研究的选择之一\scite{nlp-llm-gkd}。

这种循环改变了监督出现的位置，不能保证教师能够修复任意前缀。若学生已删除关键条件，教师看到的上下文也可能不足；若学生极不稳定，大量前缀可能没有学习价值。需要合适的学生起点、前缀筛选以及固定与在线材料的配合。Qwen3报告中的小模型蒸馏路径，也区分了较大模型提供的离线材料与学生生成上的在线监督\scite{nlp-llm-qwen3}。

在线生成表示训练数据随学生变化。只要更新仍以教师分布为监督，它就不是因为“在线”二字而变成奖励最大化。第\ref{chap:nlp-preference-alignment}章将在回答比较或奖励目标下讨论另一种更新；两种信号可以组合，但应分别写出各自改变参数的依据。

\section{微调检查与能力保持}
\label{sec:nlp-sft-checking}

训练结束只说明得到了新的参数。本节利用没有参与更新的材料，检查交互能力是否改善、原有能力是否保持，并把失败返回到相应学习环节。第\ref{chap:nlp-evaluation}章统一讨论指标和统计比较，本节先明确需要观察的对象。

\subsection{未见指令与多轮泛化}
\label{sec:nlp-sft-generalization}

若训练要求是“提取活动地点”，检查时仅换成“地点是什么”，主要检查近义表达。进一步的泛化还包括新通知、新条件组合，以及把地点与外部报名条件合并成一个回答。应分别记录这些变化，因为一个模型可能会改写模板，却不能在多条要求同时出现时完整遵守约束。

多轮检查还需要让后续问题真正依赖前文。例如先要求“以下都只依据这份通知”，再在第二轮询问材料没有给出的信息，观察模型是否说明不足；或者先要求简短回答，再增加必须包含的字段，观察模型是否在兼顾格式时遗漏内容。正确性、约束完成和历史保持应分别记录，不能把回答流畅当作完成要求的证据。

同一输入下的输出还受生成设置影响。因此，比较微调前后的模型时，需要固定输入模板、采样规则和输出预算，或者按明确协议分别检查多次生成。否则，较长预算或较低温度造成的变化可能被误认作训练收益。

\subsection{原有能力与行为边界保持}
\label{sec:nlp-sft-retention}

目标交互改善以后，还要回查原有任务。若模型学会按字段输出，却在普通问答中不断添加无关字段，格式示范已经影响了不需要该格式的场景。保持检查可以使用原有知识问答、阅读、语言与代码任务，但必须固定问题、可见信息和判断依据，才能把差异理解为模型更新后的变化。

适当拒绝也需要与正常帮助一起检查。对于材料没有提供的个人信息，说明无法确定是合适的；对于通知中明确给出的报名时间，同样拒绝则损害了正常帮助。因而，边界材料需要包括真正受限的请求、表面相似但允许的请求，以及信息不足时能够通过澄清继续完成的请求。单独提高拒绝比例不能说明模型更合适。

软蒸馏和低秩适配同样可能改变这些行为。教师可能把不必要的拒绝或固定口头禅传给学生，受限适配也可能通过共享层影响多个任务。保持要求应当参与数据与参数选择，而不是在训练结束后只附上一句“其他能力不变”。

\subsection{失败归因与下一轮微调}
\label{sec:nlp-sft-diagnosis}

对一次失败，应保留输入、可见历史、实际输出、生成设置和判定依据，再提出能够区分原因的检查。模型答错报名时间，可能是示范本身有误，也可能是截断删除了通知中的日期，还可能是模型只学会复述常见活动时间。先核对这条样本的来源、编码和计分位置，能够排除许多不需要增加训练轮次的问题。

若训练样本正确且模型仍不能拟合，可以检查更新空间是否过窄、学习是否充分；若训练任务已能完成而留出条件失败，应检查覆盖、来源差异和模板依赖；若更新后多个原有任务同时下降，则需要比较更新强度与混合策略。这里的判断是有待对照实验确认的原因假设，不能从一条错误回答直接确定是哪一个参数出了问题。

下一轮修改应对应已有证据：缺少条件就补相应示范，编码错误就重建样本，教师偏差就更换或校验信号，更新过强就改变日程或保持材料。最终检查集合应继续保留，避免不断按同一套测试答案修订材料而把检查变成另一种训练。若剩余问题主要是多种合理回答之间的取舍，单纯增加一份标准答案可能不够，下一章转向比较反馈。

\section{本章小结}
\label{sec:nlp-sft-summary}

有监督微调把期望交互行为转化为可学习的条件概率，但这条转换包含多项决定。示范决定需要模仿什么以及允许使用什么信息；角色、历史和损失掩码决定在哪里学习；教师分布决定保留哪些可接受延续；更新空间、取样和日程决定哪些变化能够发生。只有在未见指令和保持任务上检查这些变化，才能判断微调是否完成了目标。

通过这些环节，可以训练并检查模型按示范方式响应的能力；结果仍会受到示范覆盖、教师错误和参数能力的限制。第\ref{chap:nlp-preference-alignment}章将研究当回答没有唯一标准文本、但可以比较其取舍时怎样继续学习。推理轨迹也可以作为本章的示范或蒸馏材料；轨迹怎样构造、验证，以及它们是否带来推理能力，留在第\ref{chap:nlp-reasoning}章综合讨论。

\paragraph{延展阅读}
InstructGPT展示了人工示范在指令跟随流程中的位置，Self-Instruct讨论模型生成指令材料，
LIMA考察较少但经筛选的示范，LoRA则改变参数更新空间。对照这些工作时，应分别检查
示范来源、监督位置、可训练参数和独立测试条件；材料更多、参数更新更少或训练损失更低，
都不能单独证明未见任务上的行为已经改善
\scite{ouyang2022instructgpt,wang2023selfinstruct,nlp-llm-lima,ket-hu2021lora}
